\twocolumn[{%
\begin{@twocolumnfalse}
\begin{minipage}[c]{0.28\textwidth}
  \centering\includegraphics[width=\linewidth]{Imágenes/logo-udistrital.png}
\end{minipage}\hfill
\begin{minipage}[c]{0.7\textwidth}
  \centering
  {\bfseries\fontsize{17.5}{20.5}\selectfont Análisis descriptivo de propiedades fisicoquímicas y reactividad de diez hidrocarburos\par}
  \vspace{4pt}
  {\normalsize\itshape Proyecto de Química Orgánica: alcanos y alquenos\par}
  \vspace{7pt}
  {\small\bfseries Nicole Stephany Cifuentes Pardo, Maria Valentina Patiño Cadena\\
  Valeria Florez Daza, Zoreidy Taliana Rodríguez Garay, Danna Marcela Moreno Suárez\par}
  \vspace{5pt}
  {\footnotesize\itshape Universidad Distrital Francisco José de Caldas\\
  Curso: Química Orgánica \textbar\ Noviembre de 2026\par}
\end{minipage}
\vspace{9pt}
\hrule
\vspace{7pt}
\begin{minipage}{\textwidth}
\small\justifying
\noindent\textbf{Resumen.} Se analizaron diez hidrocarburos (cinco alcanos y cinco alquenos) mediante una base recopilada a partir de CAS-SciFinder. Se estudiaron diez fichas de propiedades y 200 registros de reactividad: 100 reacciones y 100 rutas de obtención. La auditoría de los archivos crudo y limpio identificó 195 números CAS de reacción distintos y concordancia en 1.060 conversiones numéricas verificables. La categoría normalizada más frecuente fue oxidación (35 registros; 17,5\%), mientras que el medio clasificado como neutro representó 38,5\%. Se dispone de temperatura numérica en 170 registros. Entre los compuestos base, butano e isobutano tienen la misma masa molecular (58,12 g/mol) y difieren 11,2~$^\circ$C en sus puntos de ebullición reportados. Las conclusiones se limitan a la muestra y consideran la heterogeneidad de condiciones, los datos faltantes y las clasificaciones automáticas.\par
\vspace{4pt}
\noindent\textbf{Palabras clave:} alcanos, alquenos, CAS-SciFinder, propiedades fisicoquímicas, reactividad, auditoría de datos.\par
\vspace{8pt}
\noindent\textbf{Abstract.} This study examines ten hydrocarbons (five alkanes and five alkenes) using records compiled from CAS SciFinder. Ten reference-compound property profiles and 200 reactivity records were analyzed: 100 reactions and 100 preparation routes. An audit of the original and cleaned workbooks identified 195 distinct CAS reaction numbers and agreement in 1,060 automatically verifiable numeric conversions. Oxidation was the most frequent normalized reaction class (35 records; 17.5\%), while the category classified as neutral accounted for 38.5\% of reported reaction media. Numerical reaction temperatures were available for 170 records. Butane and isobutane shared the same molecular mass (58.12 g/mol) but differed by 11.2~$^\circ$C in the reported boiling points. Findings are restricted to the selected records and acknowledge missing data, heterogeneous conditions, and text-based classifications.\par
\vspace{4pt}
\noindent\textbf{Keywords:} alkanes, alkenes, CAS SciFinder, physicochemical properties, reactivity, data auditing.\par
\end{minipage}
\vspace{10pt}\hrule
\vspace{17pt}
\end{@twocolumnfalse}
}]
